% Part: first-order-logic
% Chapter: beyond
% Section: second-order-logic

\documentclass[../../../include/open-logic-section]{subfiles}

\begin{document}

\olfileid{fol}{byd}{sol}

\olsection{منطق مرتبهٔ دوم}

زبان منطق مرتبهٔ دوم به ما اجازه می‌دهد نه‌تنها روی !!a{domain} از
افراد، بلکه روی رابطه‌های آن !!{domain} نیز کمیت‌گذاری کنیم. به ازای
یک زبان مرتبهٔ اول~$\Lang{L}$، برای هر $k$، !!{variable}s $R$ را
می‌افزاییم که روی رابطه‌های $k$موضعی دامنه می‌گیرند، و کمیت‌گذاری روی
این متغیرها را مجاز می‌شماریم. اگر $R$ !!a{variable} برای یک رابطهٔ
$k$موضعی باشد و $t_1$، \dots،~$t_k$ ترم‌های معمولی (مرتبهٔ اول)
باشند، $\Atom{R}{t_1,\dots,t_k}$ یک !!{formula} اتمی است. در دیگر
موارد، مجموعهٔ !!{formula}s درست مانند منطق مرتبهٔ اول، با بندهایی
اضافی برای کمیت‌گذاری مرتبهٔ دوم، تعریف می‌شود. توجه کنید که
!!{identity} را تنها برای ترم‌های مرتبهٔ اول داریم: اگر $R$ و~$S$
!!{variable}s رابطه‌ای با موضع یکسان~$k$ باشند، می‌توانیم
$\eq[R][S]$ را اختصاری برای
\[
\lforall[x_1, \dots, x_k][(\Atom{R}{x_1, \dots, x_k} \liff
  \Atom{S}{x_1, \dots, x_k})].
\]
تعریف کنیم.

قواعد منطق مرتبهٔ دوم صرفاً قواعد کمیت‌گذار را به متغیرهای جدید
مرتبهٔ دوم گسترش می‌دهند. بااین‌حال، در اینجا باید اندکی دقت کنیم و
توضیح دهیم که این متغیرها چگونه با !!{predicate}s~$\Lang{L}$ و، به
طور کلی‌تر، با !!{formula}s~$\Lang{L}$ تعامل دارند. دست‌کم، متغیرهای
رابطه‌ای ترم به شمار می‌آیند، پس استنتاج‌هایی به صورت
\[
!A(R) \Proves \lexists[R][!A(R)]
\]
داریم. اما اگر $\Lang{L}$ زبان حساب با نماد رابطه‌ای ثابت~$<$ باشد،
انتظار داریم استنتاج زیر نیز معتبر باشد:
\[
x < y \Proves \lexists[R][\Atom{R}{x,y}]
\]
یا، برای !!{formula} معین~$!A$،
\[
!A(x_1, \dots, x_k) \Proves \lexists[R][\Atom{R}{x_1,\dots,x_k}]
\]
به‌طور کلی‌تر، شاید بخواهیم استنتاج‌هایی به صورت
\[
\Subst{!A}{\lambd[\vec x][!B(\vec x)]}{R} \Proves
\lexists[R][!A]
\]
را مجاز بدانیم، که در آن
$\Subst{!A}{\lambd[\vec x][!B(\vec x)]}{R}$ حاصل جایگزین‌کردن هر
!!{formula} اتمی به صورت $\Atom{R}{t_1,\dots,t_k}$ در~$!A$ با
$!B(t_1, \dots, t_k)$ است. این قاعدهٔ اخیر معادلِ داشتن یک
{\em طرح‌وارهٔ فراگیری}، یعنی یک اصل موضوع به صورت
\[
\lexists[R][\lforall[x_1, \dots, x_k][(!A(x_1, \dots, x_k) \liff
\Atom{R}{x_1, \dots, x_k})]],
\]
به ازای هر !!{formula}~$!A$ در زبان مرتبهٔ دوم است که در آن $R$ متغیر
آزاد نیست. (تمرین: نشان دهید اگر اجازه دهیم $R$ در~$!A$ رخ دهد، این
طرح‌واره ناسازگار خواهد بود!)

منطق‌دانان معمولاً از «اصول موضوع منطق مرتبهٔ دوم» کمینه‌ترین بسط منطق
مرتبهٔ اول به‌وسیلهٔ قواعد کمیت‌گذار مرتبهٔ دوم، همراه با طرح‌وارهٔ
فراگیری، را مراد می‌کنند. اما مطالعهٔ زیرسامانه‌های ضعیف‌تر این اصول
موضوع و قواعد غالباً جالب است. برای نمونه، توجه کنید که طرح‌وارهٔ اصل
موضوع فراگیری در کلی‌ترین صورت خود \emph{ناحملی} است: به ما اجازه
می‌دهد وجود رابطهٔ $\Atom{R}{x_1, \dots, x_k}$ را که با !!a{formula}
دارای کمیت‌گذارهای مرتبهٔ دوم «تعریف» می‌شود تصدیق کنیم؛ و این
کمیت‌گذارها روی مجموعهٔ همهٔ چنین رابطه‌هایی دامنه می‌گیرند---مجموعه‌ای
که خود $R$ را نیز دربر می‌گیرد!{} در حوالی آغاز سدهٔ بیستم، واکنشی
رایج به پارادوکس راسل آن بود که تقصیر را متوجه چنین تعریف‌هایی کنند و
در پرورش مبانی ریاضیات از آن‌ها بپرهیزند. اگر کاربرد کمیت‌گذارهای
مرتبهٔ دوم را در !!{formula}~$!A$ ممنوع کنیم، صورتی \emph{حملی} از
فراگیری خواهیم داشت که قدری ضعیف‌تر است.

از دیدگاه معناشناختی، می‌توان !!{structure} مرتبهٔ دوم را متشکل از یک
!!{structure} مرتبهٔ اول برای زبان، همراه با مجموعه‌ای از رابطه‌های روی
!!{domain} دانست که کمیت‌گذارهای مرتبهٔ دوم روی آن دامنه می‌گیرند
(به‌طور دقیق‌تر، برای هر $k$ مجموعه‌ای از رابطه‌های $k$موضعی وجود
دارد). البته اگر فراگیری در دستگاه !!{derivation} گنجانده شده باشد، شرط
اضافی این است که در «بخش مرتبهٔ دوم» به‌اندازهٔ کافی رابطه وجود داشته
باشد تا اصول موضوع فراگیری را ارضا کند---وگرنه دستگاه !!{derivation}
سالم نیست!{} یک راه آسان برای تضمین وجود رابطه‌های کافی این است که بخش
مرتبهٔ دوم را متشکل از \emph{همهٔ} رابطه‌های روی بخش مرتبهٔ اول بگیریم.
چنین !!a{structure}ای \emph{تام} نامیده می‌شود و، به یک معنا، همان
«!!{structure} مقصود» واقعی برای زبان است. اگر توجه خود را به
!!{structure}s تام محدود کنیم، چیزی خواهیم داشت که \emph{معناشناسی
تام مرتبهٔ دوم} نامیده می‌شود. در این صورت، تعیین یک !!{structure} به
تعیین بخش مرتبهٔ اول فروکاسته می‌شود، زیرا محتوای بخش مرتبهٔ دوم
به‌طور ضمنی از آن به دست می‌آید.

در جمع‌بندی، سخن‌گفتن از منطق مرتبهٔ دوم قدری ابهام دارد. از حیث دستگاه
!!{derivation}، ممکن است یکی از این دو مراد باشد:
\begin{enumerate}
\item یک دستگاه !!{derivation} مرتبهٔ دوم «کمینه»، همراه با برخی اصول
  موضوع فراگیری.
\item دستگاه !!{derivation} مرتبهٔ دوم «استاندارد»، با فراگیری تام.
\end{enumerate}
از حیث معناشناسی، ممکن است به یکی از این دو علاقه‌مند باشیم:
\begin{enumerate}
\item معناشناسی «ضعیف»، که در آن !!a{structure} از یک بخش مرتبهٔ اول
  همراه با بخشی مرتبهٔ دوم که برای ارضای اصول موضوع فراگیری به‌اندازهٔ
  کافی بزرگ است تشکیل می‌شود.
\item معناشناسی «استاندارد» مرتبهٔ دوم، که در آن تنها
  !!{structure}s تام را در نظر می‌گیریم.
\end{enumerate}
هنگامی که منطق‌دانان دستگاه !!{derivation} یا معناشناسی مورد نظر خود را
مشخص نمی‌کنند، معمولاً مورد دوم هر فهرست را مراد می‌کنند. مزیت کاربرد
این معناشناسی آن است که، چنان‌که خواهیم دید، توصیف‌هایی مقوله‌ای از
بسیاری ساختارهای طبیعی ریاضی به دست می‌دهد؛ هم‌زمان، دستگاه
!!{derivation} بسیار قوی و برای این معناشناسی سالم است. عیب آن این است
که دستگاه !!{derivation} برای معناشناسی \emph{تمام} نیست؛ در واقع،
\emph{هیچ} دستگاه !!{derivation}ای که به‌نحو مؤثر داده شده باشد برای
معناشناسی تام مرتبهٔ دوم تمام نیست. از سوی دیگر، خواهیم دید که دستگاه
!!{derivation} برای معناشناسی ضعیف‌شده \emph{تمام} است؛ از این نتیجه
می‌شود که اگر جمله‌ای اثبات‌پذیر نباشد، \emph{یک} !!{structure}، نه
لزوماً ساختار تام، وجود دارد که آن جمله در آن کاذب است.

زبان منطق مرتبهٔ دوم بسیار غنی است. می‌توان رابطه‌های یک‌موضعی را با
زیرمجموعه‌های !!{domain} یکی گرفت و، در نتیجه، به‌ویژه روی این مجموعه‌ها
کمیت‌گذاری کرد؛ برای نمونه، می‌توان استقرا برای اعداد طبیعی را با یک
اصل موضوع بیان کرد:
\[
\lforall[R][((\Atom{R}{\Obj{0}} \land \lforall[x][(\Atom{R}{x} \lif
    \Atom{R}{x'})]) \lif \lforall[x][\Atom{R}{x}])].
\]
اگر زبان حساب را دارای نمادهای $\Obj 0, \Obj \prime, +,
\times$ و~$<$ بگیریم، می‌توانیم اصول موضوع زیر را برای توصیف رفتار
آن‌ها بیفزاییم:
\begin{enumerate}
\item $\lforall[x][\lnot x' = \Obj 0]$
\item $\lforall[x][\lforall[y][(x' = y' \lif x = y)]]$
\item $\lforall[x][(x + \Obj 0) = x]$
\item $\lforall[x][\lforall[y][(x + y') = (x + y)']]$
\item $\lforall[x][(x \times \Obj 0) = \Obj 0]$
\item $\lforall[x][\lforall[y][(x \times y') = ((x \times y) + x)]]$
\item $\lforall[x][\lforall[y][(x < y \liff \lexists[z][y = (x + z')])]]$
\end{enumerate}
نشان‌دادن اینکه این اصول موضوع، همراه با اصل موضوع استقرای بالا، توصیفی
مقوله‌ای از !!{structure}~$\Struct{N}$، مدل استاندارد حساب، به دست
می‌دهند دشوار نیست، مشروط به اینکه از معناشناسی تام مرتبهٔ دوم استفاده
کنیم. هر !!{structure}~$\Struct{M}$ را که این اصول موضوع در آن صادق‌اند
در نظر بگیرید، و با بازگشت معمول روی~$\Nat$ تابعی~$f$ از~$\Nat$ به
!!{domain} ساختار~$\Struct{M}$ تعریف کنید، به‌گونه‌ای که
$f(0) = \Assign{\Obj 0}{M}$ و
$f(x+1) = \Assign{\prime}{M}(f(x))$. با استقرای معمول روی~$\Nat$ و
این واقعیت که اصول موضوع (1) و~(2) در~$\Struct M$ برقرارند، می‌بینیم
که $f$ !!{injective} است. برای دیدن اینکه $f$ !!{surjective} است،
مجموعهٔ عناصر~$\Domain{M}$ را که در برد~$f$ هستند با~$P$ نشان دهید.
چون $\Struct M$ تام است، $P$ در !!{domain} مرتبهٔ دوم قرار دارد. بنا بر
ساخت~$f$، می‌دانیم $\Assign{\Obj 0}{M}$ در~$P$ است و $P$ تحت
$\Assign{\prime}{M}$ بسته است. این واقعیت که اصل موضوع استقرا در
$\Struct M$ (به‌ویژه برای~$P$) برقرار است تضمین می‌کند که $P$ با تمام
!!{domain} مرتبهٔ اول~$\Struct M$ برابر است. از اینجا معلوم می‌شود که
$f$ !!a{bijection} است. نشان‌دادن اینکه $f$ یک همریختی است نیز دشوارتر
نیست و با کاربرد مکرر استقرای معمول روی~$\Nat$ انجام می‌شود.

در بیان نظریهٔ مجموعه‌ای، تابع فقط نوع خاصی از رابطه است؛ برای نمونه،
یک تابع یک‌موضعی~$f$ را می‌توان با رابطه‌ای دوموضعی~$R$ که
$\lforall[x][\lexists![y][R(x,y)]]$ را برآورده می‌کند یکی گرفت. در
نتیجه، می‌توان روی تابع‌ها نیز کمیت‌گذاری کرد. سپس با استفاده از
معناشناسی تام می‌توان ردهٔ !!{structure}s نامتناهی را ردهٔ آن
!!{structure}s $\Struct M$ تعریف کرد که برایشان تابعی یک‌به‌یک از
!!{domain}~$\Struct M$ به یک زیرمجموعهٔ سرهٔ خودش وجود دارد:
\[
\lexists[f][(\lforall[x][\lforall[y][(\eq[f(x)][f(y)] \lif \eq[x][y])]]
  \land \lexists[y][\lforall[x][\eq/[f(x)][y]]])].
\]
نقیض این جمله آنگاه ردهٔ !!{structure}s متناهی را تعریف می‌کند.

افزون بر این، می‌توان با افزودن عبارت زیر به تعریف یک ترتیب خطی، ردهٔ
خوش‌ترتیب‌ها را تعریف کرد:
\[
\lforall[P][(\lexists[x][\Atom{P}{x}] \lif \lexists[x][(\Atom{P}{x}
    \land \lforall[y][(y < x \lif \lnot \Atom{P}{y})])])].
\]
این عبارت، با یکی‌گرفتن «مجموعه» و «رابطهٔ یک‌موضعی»، می‌گوید هر
مجموعهٔ ناتهی دارای کوچک‌ترین عنصر است. به‌عنوان نمونه‌ای دیگر، می‌توان
مفهوم همبندی برای گراف‌ها را با گفتن اینکه هیچ جداسازی نابدیهی رأس‌ها
به بخش‌های ناهمبند وجود ندارد بیان کرد:
\[
\lnot \lexists[A][(\lexists[x][A(x)] \land \lexists[y][\lnot A(y)]
  \land \lforall[w][\lforall[z][((\Atom{A}{w} \land \lnot \Atom{A}{z})
      \lif \lnot \Atom{R}{w,z})]])].
\]
به‌عنوان نمونه‌ای دیگر، می‌توانید در قالب تمرین بکوشید ردهٔ
!!{structure}s متناهی‌ای را تعریف کنید که !!{domain} آن‌ها اندازه‌ای
زوج دارد. چشمگیرتر آنکه می‌توان توصیفی مقوله‌ای از اعداد حقیقی به‌مثابه
یک میدان مرتب کامل شامل اعداد گویا به دست داد.

خلاصه اینکه منطق مرتبهٔ دوم بسیار بیانگرتر از منطق مرتبهٔ اول است. این
خبر خوب بود؛ اکنون خبر بد. پیش‌تر گفتیم که هیچ دستگاه !!{derivation}
مؤثری وجود ندارد که برای معناشناسی تام مرتبهٔ دوم تمام باشد. خواه خوش
آیند خواه ناخوش، بسیاری از ویژگی‌های منطق مرتبهٔ اول غایب‌اند، از جمله
فشردگی و قضایای لوونهایم--اسکولم.

از سوی دیگر، اگر آماده باشیم معناشناسی تام مرتبهٔ دوم را به سود
معناشناسی ضعیف‌تر کنار بگذاریم، آنگاه دستگاه !!{derivation} کمینهٔ
مرتبهٔ دوم برای این معناشناسی تمام است. به بیان دیگر، اگر $\Proves$ را
«اثبات می‌کند در دستگاه کمینه» و $\Entails$ را «در معناشناسی ضعیف‌تر
منطقاً نتیجه می‌دهد» بخوانیم، می‌توانیم نشان دهیم هرگاه
$\Gamma \Entails !A$، آنگاه $\Gamma \Proves !A$. اگر بخواهیم اصول
موضوع فراگیری خاصی را در دستگاه !!{derivation} بگنجانیم، باید معناشناسی
را به ساختارهای مرتبهٔ دومی محدود کنیم که این اصول موضوع را ارضا
می‌کنند: برای نمونه، اگر $\Delta$ از مجموعه‌ای از اصول موضوع فراگیری
(احتمالاً همهٔ آن‌ها) تشکیل شده باشد، هرگاه
$\Gamma \cup \Delta \Entails !A$، آنگاه
$\Gamma \cup \Delta \Proves !A$. به‌ویژه، اگر $!A$ با استفاده از اصول
موضوع فراگیری مورد نظرمان اثبات‌پذیر نباشد، آنگاه مدلی برای~$\lnot !A$
وجود دارد که بااین‌همه این اصول موضوع فراگیری در آن برقرارند.

ساده‌ترین راه برای دیدن اینکه قضیهٔ تمامیت برای معناشناسی ضعیف‌تر
برقرار است، آن است که منطق مرتبهٔ دوم را به صورت منطق چندنوعی زیر در
نظر بگیریم. یک نوع به !!{domain} معمول «مرتبهٔ اول» تعبیر می‌شود و سپس
برای هر $k$ !!a{domain} از «رابطه‌های $k$موضعی» داریم. زبان را دارای
نمادهای رابطه‌ای درونیِ «$\Atom{\Obj{true}_k}{R,x_1,\dots,x_k}$»
می‌گیریم که مقصود از آن‌ها این است که $R$ دربارهٔ $x_1$، \dots،~$x_k$
برقرار است، جایی که $R$ متغیری از نوع «رابطهٔ $k$موضعی» و
$x_1$، \dots،~$x_k$ اشیایی از نوع مرتبهٔ اول‌اند.

با این یکی‌گرفتن، معناشناسی ضعیف مرتبهٔ دوم اساساً همان معناشناسی معمول
منطق چندنوعی است؛ و پیش‌تر دیدیم که منطق چندنوعی را می‌توان در منطق
مرتبهٔ اول نشاند. بنابراین، با صرف نظر از ترجمه‌های رفت‌وبرگشتی، تلقی
ضعیف‌تر از منطق مرتبهٔ دوم در واقع صورتی از منطق مرتبهٔ اول در لباس
مبدل است که در آن !!{domain} شامل هم «اشیا» و هم «رابطه‌ها»ست و اصول
موضوع مناسب بر آن‌ها حاکم‌اند.

\end{document}
